\documentclass[10pt]{letter}
\usepackage[letterpaper,margin=0.85in]{geometry}
\usepackage[T1]{fontenc}
\usepackage[utf8]{inputenc}
\usepackage{microtype}
\usepackage[colorlinks=true,urlcolor=blue]{hyperref}


\begin{document}
\begin{letter}{Dr.\ Chun-Wei Yang\\Handling Editor\\Scientific Reports}
\opening{Dear Dr.\ Yang,}

Please find the major revision of \emph{RankCloak: Hiding Synthetic Cryptographic Artifacts in Plain English with Language Models} (submission 0170c7b2-1065-4344-a5d5-b839b51ac22c). I thank you and the reviewers for the detailed guidance. The manuscript replaces the original evidence with the more thorough study and is accompanied by a point-by-point response.

The revision expands the corpus to 480 public deterministic outputs from eight declared cryptographic classes, evaluated across three open-source model families and 18 English prompt templates in six categories. It adds grouped inference, controlled transmission perturbations and first-divergence analysis, automated readability and held-out quality diagnostics, filter/lead-in/segment/tail ablations, prompt-topic contrasts, formal capacity relationships, secondary Spanish and the 'Simplified Chinese' exact-copy recovery, and computational-overhead measurements.

I also added a 15,840-row neural-steganalysis study using TextCNN-equivalent and DeBERTa-v3-base classifiers across matched and held-out regimes, with leakage and near-duplicate sensitivity audits. Matched ROC-AUC was near one for both architectures and is presented as adverse evidence. The revised conclusions do not claim undetectability, secure or robust communication, or deployment readiness.

The recovery claim is now narrow with all 6,480 primary serialized payloads recovered under saved-token exact-copy replay, but visible-text retokenization and common transformations were fragile. Automated measures are not described as human ratings; no participant study or deployment was performed. The manuscript discloses unavailable cells, singular fits, post-outcome analyses, missing local GGUF files, incomplete perturbation coverage, and timing-scope limits.

The combined Data and Code availability section points to the public repository and the published Zenodo code-and-data archive DOI, \url{https://doi.org/10.5281/zenodo.21987450}.

I believe the revision preserves RankCloak's methodological contribution while giving the adverse recoverability and detectability evidence it needs. Thank you for considering the revised manuscript.

Sincerely,

Alexander V. Mantzaris, Associate Professor \\
School Data, Mathematical, and Statistical Sciences, UCF \\
4000 University Bvld, Orlando FL, 32817 USA \\
alexander.mantzaris@ucf.edu \\
TCII 211C, +1-407-692-0280
\end{letter}
\end{document}
